\documentclass[12pt, a4paper]{article}

% --- Packages ---
\usepackage[utf8]{inputenc}
\usepackage[T1]{fontenc}
\usepackage{lmodern}
\usepackage[english]{babel}
\usepackage{amsmath, amssymb}
\usepackage{graphicx}
\usepackage[hidelinks]{hyperref}
\usepackage{booktabs}
\usepackage{caption}
\usepackage{subcaption}
\usepackage[margin=2.5cm]{geometry}
\usepackage{setspace}
\usepackage{natbib}
\usepackage{float}
\usepackage{enumitem}
\usepackage{xcolor}
\usepackage{array}

\onehalfspacing

% --- Title ---
\title{\textbf{Digital Mirrors in Motion: Measuring Gendered Movement in Short Videos}\thanks{I thank Tsolak and K\"uhne for making their dataset publicly available. Full source code and implementation details are available at \url{https://github.com/FlyinHedgehog/unveiling-digital-mirrors-further-research}.}}
\author{G\"ok\c{c}e \c{S}ahin\thanks{LMU Munich, MSc Data Science}}
\date{March 15, 2026\\[0.5em]Course: Computational Social Science (CSS)}

\begin{document}
\maketitle

% --- Abstract ---
\begin{abstract}
Social media platforms have evolved rapidly from static imagery to dynamic short video content, yet sociological methods for analyzing gendered self-representation remain largely rooted in static image analysis. Building on the recent work of \citet{tsolak2025unveiling}, who identified 150 gendered body-posing clusters in Instagram imagery, this paper introduces a computational tool designed to extend this methodology to video data. Using the original raw keypoint dataset, I independently derive a set of 150 pose prototypes through K-means clustering. These prototypes form the foundation of an open-source web application \citep{sahin2026unveiling} that uses MediaPipe for pose estimation. To address the domain gap between datasets, the tool employs a custom mapping layer to align live video landmarks with the training clusters. This framework enables scalable extraction and analysis of gendered body patterns in dynamic digital environments, helping bridge the gap between sociological theory and modern content formats.
\medskip
\noindent\textbf{Keywords:} Computational Social Science, Computer Vision, Gender Performativity, Video Analysis, Social Media
\end{abstract}

\newpage
\tableofcontents
\newpage

% =============================================================================
% I. INTRODUCTION
% =============================================================================
\section{Introduction}
\label{sec:intro}

Social media platforms have fundamentally changed how individuals perform identity, acting as ``digital mirrors'' where users share self-representations for public consumption. Theoretical frameworks regarding the presentation of self, most prominently developed by \citet{goffman1979gender}, have found new relevance in the digital age, where likes, comments, and algorithmic visibility incentivize specific forms of gender performance. Recent sociological research has begun to decode these representations quantitatively; notably, \citet{tsolak2025unveiling} employed OpenPose \citep{openpose2019} to identify 150 distinct body-posing prototypes from a massive dataset of Instagram images. Their findings highlighted that some of these prototypes showed significant gender differences, providing empirical evidence for the persistence of gender stereotypes in digital mirrors.

However, digital platforms are evolving rapidly. The dominance of image feeds is long gone, and users are now surrounded by short video contents driven by platforms such as TikTok, YouTube Shorts, and Instagram Reels. While \citet{tsolak2025unveiling} provided a scalable methodology for full-body pose images, current social science research lacks accessible, scalable tools to apply these rigorous clustering methods to dynamic video data. Recent work by \citet{rittmann2024body} demonstrates how pose estimation can quantify nonverbal behavior in video, specifically measuring gesticulation and posture in political speeches, but this approach has not yet been applied to the domain of gendered self-representation.

This paper introduces a computational social science tool \citep{sahin2026unveiling} designed to bridge this gap. Using the labeled dataset provided by \citet{tsolak2025unveiling}, I implemented an independent unsupervised learning pipeline to derive a set of 150 body-pose prototypes. These prototypes serve as the analytical foundation for a web-based application capable of extracting, normalizing, and classifying body poses from video streams on a frame-by-frame basis. To enable real-time analysis in the browser, I use the lightweight MediaPipe \citep{mediapipe2019} framework and a custom mapping layer to align uploaded video data with the original OpenPose-derived clusters. This framework allows researchers to move beyond static, single-frame analysis and observe the temporal dynamics of gendered body stereotyping as they unfold in real time.


% =============================================================================
% II. BACKGROUND
% =============================================================================
\section{Background}
\label{sec:background}

\subsection{Sociological Context}
\label{sec:sociology}

The sociological analysis of gendered body language is grounded in the concept of ``doing gender'' \citep{west1987doing}, which posits that gender is not a fixed trait but an accomplished activity that is continuously performed in interaction. \citet{goffman1979gender}'s earlier analysis of gender advertisements identified specific codes, such as the ``feminine touch,'' that signal gender identity through posture and gesture. In the context of social media, these performances are directed not only at human audiences but are also mediated by algorithmic feedback loops that may reinforce stereotypical displays.

Building on Goffman, more recent research has found that these gendered display patterns persist in digital contexts. \citet{doring2016selfies} analyzed 500 selfies on Instagram and confirmed the presence of gender-stereotypical self-presentation modes, including the feminine touch and imbalanced postures among women, as well as muscle display among men. \citet{gotz2019bein} identified additional patterns, such as S-shaped body displays and asymmetric leg crossing among female influencers. Meanwhile, the concept of ``manspreading'' (wide-legged sitting postures predominantly adopted by men) has been documented across both offline and online settings \citep{jane2017manspreading}.

\citet{butler1990gender} advanced the theoretical understanding of gender by arguing that gender identity is constituted through the repeated ``stylization of the body,'' thereby framing it as performative rather than essential. This perspective is especially relevant to social media, where users actively curate their self-presentation for audiences. Taken together, these theoretical and empirical foundations suggest that gendered body posing is not incidental but constitutive of gender identity and leaves traceable patterns in visual data.


\subsection{The Digital Mirrors Dataset}
\label{sec:dataset}

This research builds directly upon the empirical foundation laid by \citet{tsolak2025unveiling}. In their study, they analyzed 832,667 images collected from the Instagram API between June 2017 and October 2018, utilizing the OpenPose library \citep{openpose2019} to extract 2D skeletal data from detected individuals. After filtering for single-person, full-body shots with high-confidence keypoint detection and crowdsourcing gender annotations via Amazon Mechanical Turk, the authors arrived at a final analysis sample of 15,167 images (52.46\% women, 43.50\% men, 4.04\% non-binary persons).

Using an unsupervised learning approach that combines hierarchical clustering with K-means ($k=150$), they identified 150 unique clusters of body poses. Each cluster yields a \textit{prototype}, defined as the arithmetic mean of all member pose vectors, which can be visualized as a stick-figure summary of the cluster's characteristic posture. Statistical analysis using multinomial tests with Bonferroni correction ($\alpha = 0.001/150$) revealed that approximately 20\% of these clusters exhibited a statistically significant gender imbalance, providing empirical evidence for the persistence of gendered body posing in social media imagery.

Their analysis revealed that while many poses are gender-neutral, specific clusters are strongly coded as masculine or feminine. Wide stances and ``manspreading'' were statistically overrepresented in images of men, while asymmetric postures, crossed legs, S-shaped body displays, and the ``feminine touch'' were predominant in images of women.

% =============================================================================
% III. METHOD
% =============================================================================

\section{Method}
\label{sec:method}

\subsection{System Architecture}
\label{sec:architecture}

The proposed framework is implemented as a client-server web application. The backend consists of Python scripts used for initial data processing and clustering of the training dataset. The frontend is a JavaScript-based interface that handles real-time user interaction, pose estimation, and classification directly in the browser. This architecture ensures privacy and efficiency, as video data are processed client-side without requiring upload to an external server. 

The tool offers three analysis modes: (1)~a \textbf{single-image} mode that performs pose detection and cluster matching on uploaded images, (2)~a \textbf{video analysis} mode that extracts poses with a 0.5 second intervals and generates a temporal profile of gendered posing, and (3)~a \textbf{pose database} browser that allows researchers to explore all 150 cluster prototypes, filter by gender category, and inspect the underlying gender distributions.

Figure~\ref{fig:pipeline} in the Appendix provides a comprehensive overview of the computational framework, detailing how the foundational dataset was processed into labeled prototypes alongside the application logic used to classify new video and image inputs.


\subsection{Data and Pose Estimation}
\label{sec:pose_estimation}

The analytical foundation of this tool is the dataset provided by \citet{tsolak2025unveiling}, which comprises 15,167 annotated images. The dataset uses a 17-keypoint skeletal structure, including anatomical landmarks such as the ``Neck'' and ``MidHip.''

To enable real-time video analysis, our tool utilizes MediaPipe Pose \citep{mediapipe2019}, which natively infers 33 3D landmarks. A critical methodological challenge was the domain gap between MediaPipe's output and the 17-point structure of the training data. To address this, I implemented a mapping layer that discards the depth (Z) coordinate and translates the remaining 2D points to the 17-point format.

\begin{itemize}[nosep]
    \item \textbf{Direct Mapping:} 15 landmarks (Nose, Shoulders, Elbows, Wrists, Hips, Knees, Ankles) are mapped directly from MediaPipe to the target format. Additionally, MediaPipe's ``left/right foot index'' landmarks are mapped to the ``Big Toe'' keypoints from the training data, as these are the closest anatomical equivalents.
    \item \textbf{Derived Landmarks:} MediaPipe does not provide ``Neck'' or ``MidHip'' keypoints. I derive these geometrically: the ``Neck'' is calculated as the midpoint between the left and right shoulders, and the ``MidHip'' as the midpoint between the left and right hips.
\end{itemize}

\noindent Landmarks with a visibility score below 0.5 are filtered out to ensure data quality. 


\subsection{Normalization Pipeline}
\label{sec:normalization}

To ensure scale and translation invariance, I implemented the exact normalization logic described by \citet{tsolak2025unveiling}. For a set of keypoints $P = \{(x_i, y_i)\}_{i=1}^{n}$, the bounding box is determined by the minimum and maximum $x$ and $y$ coordinates.

First, the pose is translated such that the bottom-left corner of the bounding box aligns with the origin~$(0, 0)$. Note that in image coordinate systems (where $y$ increases downward), the bottom-left is defined by $(x_{\min}, y_{\max})$:

\begin{equation}
    x'_i = x_i - x_{\min}, \quad y'_i = y_i - y_{\max}
    \label{eq:translation}
\end{equation}

Second, the translated coordinates are flattened into a vector $\mathbf{v}$ and normalized to unit length using the $L_2$ norm:

\begin{equation}
    \mathbf{v}_{\text{normalized}} = \frac{\mathbf{v}}{\|\mathbf{v}\|_2}
    \label{eq:normalization}
\end{equation}

\noindent This two-step normalization ensures that distance computations between new poses and cluster centroids are invariant to the scale and position of the detected person within the frame.


\subsection{Clustering and Labeling Strategy}
\label{sec:clustering}

Using the normalized training data, I performed K-means clustering ($k = 150$) to generate a set of pose prototypes. The value of $k = 150$ was adopted directly from \citet{tsolak2025unveiling}, who determined it through a two-step procedure: agglomerative hierarchical clustering (Ward's method), followed by qualitative examination of the resulting dendrogram. Because this study clusters the same dataset, there is no need to rederive $k$; the original authors' choice remains directly applicable. Cluster prototypes were computed as the arithmetic mean of all member pose vectors.

\noindent This clustering was performed independently of the gender annotations; labels were introduced only at the statistical analysis stage, consistent with the methodology of \citet{tsolak2025unveiling}.

To assign gender labels to these independent clusters, I applied a chi-squared goodness-of-fit test against the null hypothesis derived from the dataset distribution (Females: 52.46\%, Males: 43.50\%, Non-binary: 4.04\%). Following the methodology of \citet{tsolak2025unveiling}, I adopted a strictly conservative threshold of $\alpha = 0.001$ before Bonferroni correction. The corrected critical value is thus $\alpha_{\text{corrected}} = 0.001 / 150 \approx 6.667 \times 10^{-6}$. To ensure statistical robustness and prevent artificial inflation of significance, I excluded clusters where the expected frequency in any gender category was lower than 1. This safeguard is necessary to avoid unreliable $p$-values, particularly in smaller clusters where the low (4.04\%) baseline prevalence of non-binary individuals would fail to meet the minimum frequency requirements for the chi-squared approximation. Clusters with a statistically significant over-representation of a specific gender were labeled ``Masculine'' or ``Feminine,'' while those with no significant deviation were labeled ``Neutral.''




\subsection{Video Analysis Logic}
\label{sec:video_logic}

The video analysis module extracts frames at a configurable interval (default: 0.5 seconds). Each frame is processed independently.

\begin{enumerate}[nosep]
    \item The pose is detected via MediaPipe and mapped to the 17-point format.
    \item The pose is normalized using the $L_2$ unit-vector formula described in Section~\ref{sec:normalization}.
    \item The system calculates the Euclidean distance against the entire dataset of individual normalized poses (15,167 poses) to find the single nearest neighbor, and then assigns the frame to the parent cluster category of that nearest neighbor.
\end{enumerate}

\noindent Frames where no pose is detected are skipped. The tool aggregates these frame-level classifications to generate a summary of the gendered pose performance over the video's duration, including counts and percentages of masculine, feminine, and neutral frames. A timeline chart provides a visual representation of the classification sequence, allowing researchers to observe temporal shifts in gendered posing.




% =============================================================================
% IV. RESULTS
% =============================================================================
\section{Results}
\label{sec:results}

\subsection{Clustering Output}
\label{sec:clustering_results}

All 150 clusters contained at least one observation. The median cluster size was 54 poses, with a minimum of 6 and a maximum of 437. Thirty clusters exhibited a statistically significant deviation from the expected gender distribution ($p < 6.667 \times 10^{-6}$, chi-squared test with Bonferroni correction). Among these, 15 clusters were classified as masculine (male overrepresentation) and 15 as feminine (female overrepresentation). The remaining 120 clusters showed no statistically significant gender imbalance and were classified as neutral. No statistically significant non-binary prototype was identified.

Figure~\ref{fig:all_clusters} shows the gender composition of all 150 clusters, ordered by the proportion of female-labeled poses. The gradient from male-dominated (left) to female-dominated (right) clusters illustrates the spectrum of gendered posing across the entire cluster space.

\begin{figure}[H]
    \centering
    \includegraphics[width=\textwidth]{Figures/gender_distribution_all.png}
    \caption{Gender composition of all 150 clusters, sorted by the proportion of female-labeled poses. Each bar shows the relative share of women (orange), men (teal), and non-binary persons (blue) within the cluster. The gradient illustrates the spectrum from predominantly masculine clusters (left) to predominantly feminine clusters (right).}
    \label{fig:all_clusters}
\end{figure}

Figure~\ref{fig:significant_clusters} highlights the 30 clusters that reached statistical significance after Bonferroni correction, with annotated sample sizes.

\begin{figure}[H]
    \centering
    \includegraphics[width=\textwidth]{Figures/gender_distribution_significant.png}
    \caption{Gender composition of statistically significant clusters ($p < 6.667 \times 10^{-6}$, Bonferroni-corrected). Bars are annotated with absolute counts per gender category. Orange = women, teal = men, blue = non-binary persons.}
    \label{fig:significant_clusters}
\end{figure}

Table~\ref{tab:cluster_summary} presents a summary of cluster classification outcomes.

\begin{table}[ht]
    \centering
    \caption{Summary of cluster classification by gender character.}
    \label{tab:cluster_summary}
    \begin{tabular}{lcc}
        \toprule
        \textbf{Category} & \textbf{Number of Clusters} & \textbf{Percentage} \\
        \midrule
        Masculine  & 15  & 10\%  \\
        Feminine   & 15  & 10\%  \\
        Neutral    & 20  & 80\%  \\
        \midrule
        Total      & 150   & 100\% \\
        \bottomrule
    \end{tabular}
\end{table}



\subsection{Tool Capabilities}
\label{sec:tool_capabilities}

The web application was successfully deployed as a client-side application requiring no server-side computation beyond static file hosting. During testing, the tool demonstrated the ability to:

\begin{itemize}[nosep]
    \item Detect and normalize poses from uploaded videos.
    \item Match each extracted frame to one of the 150 cluster prototypes by Euclidean distance.
    \item Generate a temporal pose timeline showing the sequence of masculine, feminine, and neutral classifications across a video's duration.
    \item Export per-frame data as CSV for statistical analysis.
\end{itemize}


The web application's primary interface for video analysis is shown in Figure~\ref{fig:tool_screenshot_video}. Upon uploading a video, the system performs temporal sampling and landmark extraction. The interface displays a gallery of extracted frames where each frame is overlaid with the detected MediaPipe skeleton. This allows researchers to visually inspect the accuracy of the pose estimation before proceeding to aggregate analysis.

\begin{figure}[H]
    \centering
    \includegraphics[width=0.78\textwidth,height=0.60\textheight,keepaspectratio]{Figures/app-ss-video.png}
    \caption{Screenshot of the video analysis interface showing the extracted frame gallery.}
    \label{fig:tool_screenshot_video}
\end{figure}

Beyond individual frame inspection, the tool provides a temporal overview of the video's gendered posing performance. As shown in Figure~\ref{fig:tool_screenshot_timeline}, the pose timeline visualizes the sequence of classifications (Masculine, Feminine, or Neutral) throughout the video's duration. This view is particularly useful for identifying specific moments of gendered transitions or sustained stereotypical performances.

\begin{figure}[H]
    \centering
    \includegraphics[width=0.78\textwidth,height=0.60\textheight,keepaspectratio]{Figures/app-ss-timeline.png}
    \caption{The pose timeline chart, visualizing the temporal dynamics of masculine and feminine poses in a video.}
    \label{fig:tool_screenshot_timeline}
\end{figure}

For each frame, the "Current Pose" panel (Figure~\ref{fig:tool_screenshot_match}) provides granular details about the classification. It displays the gender distribution (male vs. female percentage) of the matched cluster and provides a side-by-side comparison between the cluster's prototype and the specific dataset pose that constitutes the nearest neighbor.

\begin{figure}[H]
    \centering
    \includegraphics[width=0.78\textwidth,height=0.60\textheight,keepaspectratio]{Figures/app-ss-match.png}
    \caption{The match prediction panel, showing gender percentages and the nearest individual pose match.}
    \label{fig:tool_screenshot_match}
\end{figure}

For researchers interested in exploring the categorical basis of these classifications, the Database Browser (Figure~\ref{fig:tool_screenshot_database}) provides an interactive way to explore the 150 pose clusters. Users can filter clusters by their gendered character (Masculine, Feminine, or Neutral) and see aggregate statistics for each category.

\begin{figure}[H]
    \centering
    \includegraphics[width=0.78\textwidth,height=0.60\textheight,keepaspectratio]{Figures/app-ss-database.png}
    \caption{The database patterns browser, allowing users to browse pose clusters by gender category.}
    \label{fig:tool_screenshot_database}
\end{figure}

Finally, clicking on any cluster reveals its descriptive details (Figure~\ref{fig:tool_screenshot_prototype}). This view presents the cluster's average prototype alongside a selection of diverse member poses from the original dataset, enabling a qualitative understanding of the morphological variations within a single pose category.

\begin{figure}[H]
    \centering
    \includegraphics[width=0.78\textwidth,height=0.60\textheight,keepaspectratio]{Figures/app-ss-prototype.png}
    \caption{Deep-dive view of a single cluster, displaying its prototype and multiple representative member poses.}
    \label{fig:tool_screenshot_prototype}
\end{figure}



\subsection{Validation Against Original Study}
\label{sec:validation}

Because the clustering was performed independently (using scikit-learn's K-Means implementation with $k = 150$, \texttt{n\_init=20}, \texttt{max\_iter=500}, and \texttt{random\_state=34}), the resulting cluster assignments do not correspond one-to-one with the cluster IDs reported by \citet{tsolak2025unveiling}. However, qualitative comparison of the prototype stick figures (see Figures~\ref{fig:masc_prototypes}, \ref{fig:fem_prototypes}, and \ref{fig:neutral_prototypes} in the Appendix) reveals consistent recovery of the major gendered pose categories documented in the original paper.

\begin{itemize}[nosep]
    \item \textbf{Wide-stance and symmetric postures} (predominantly masculine) appear in multiple clusters, consistent with the ``power posing'' patterns.
    \item \textbf{Crossed-leg and asymmetric postures} (predominantly feminine) recur across several clusters, reflecting the ``imbalance'' and ``S-shaped body display'' themes well-documented in gender studies \citep{doring2016selfies}.
    \item \textbf{Seated, legs-apart postures} (predominantly masculine) reproduce the ``manspreading'' clusters identified by \citet{tsolak2025unveiling}.
\end{itemize}

\noindent This qualitative consistency supports the robustness of the clustering methodology and validates the use of these prototypes as a basis for video classification.


% =============================================================================
% VI. DISCUSSION
% =============================================================================

\section{Discussion}
\label{sec:discussion}

\subsection{Interpretation of Results}
\label{sec:interpretation}

The successful deployment of the web application demonstrates that the 150 pose prototypes derived from static Instagram imagery are robust enough to classify dynamic video frames. The tool effectively captures the fluidity of gender performance; preliminary tests suggest that individuals do not maintain a single ``gendered'' pose throughout a video but rather drift in and out of gender-coded clusters. This supports the sociological view of gender as an active, continuous performance rather than a fixed attribute.


\subsection{Limitations}
\label{sec:limitations}

Despite the tool's utility, several limitations should be acknowledged. First, the analysis is constrained to two-dimensional space. Video content often involves complex camera angles and foreshortening, which can distort limb lengths. Although the $L_2$ normalization strategy mitigates scale differences, it cannot fully correct for perspective distortion, potentially leading to misclassification.

Second, a domain gap remains between the training data (OpenPose) and the inference engine (MediaPipe). Although the custom mapping layer, which derives Neck and MidHip points, helps harmonize the skeletal structures, minor discrepancies in keypoint confidence and placement logic between the two models may still introduce noise.

Third, the training data originate from a specific geographic context (Berlin) and time period (2017--2018), and the gender annotations are based on perceived gender labels assigned by crowd workers rather than self-reported identity. These constraints, inherited from the original \citet{tsolak2025unveiling} dataset, limit the generalizability of the prototypes across cultural contexts and time periods. Additionally, the non-binary category in the training data is relatively small (4.04\%), which limits the tool's ability to identify poses specifically associated with non-binary individuals.

Finally, the tool currently analyzes only single-person, full-body poses. Extending the analysis to multi-person scenes, partial body shots, or close-up framing would require additional methodological development beyond the scope of this paper.


\subsection{Future Research: AI-Generated Content}
\label{sec:future_work}

A promising avenue for future research involves the analysis of AI-generated video content. As generative AI models increasingly produce realistic human movements, there is a risk that these models might inadvertently encode or even amplify existing gender-based pose stereotypes. Initial qualitative observations of AI-generated material suggest that these tools may follow highly conventionalized and exaggerated gendered body signatures. Extending the current framework to large-scale datasets of synthetic media would allow researchers to quantify whether digital mirrors constructed by algorithms are more gender-stereotypical than those created by human users.


\subsection{Ethical Considerations}
\label{sec:ethics}

Automated analysis of gender involves inherent ethical risks. It is crucial to clarify that this tool classifies \emph{poses}, not people. The labels ``Masculine'' and ``Feminine'' refer to sociological codes established in the training dataset; they describe the statistical gender composition of a cluster rather than the biological sex or gender identity of the subject. Misinterpretation of these results could lead to essentialist claims. Therefore, this tool should be used to analyze societal patterns of representation rather than to categorize individual users.


\section{Conclusion}
\label{sec:conclusion}

This paper introduced a computational tool that extends the body-pose analysis methodology of \citet{tsolak2025unveiling} from static Instagram imagery to dynamic video content. By independently replicating the clustering pipeline on the original dataset and building a browser-based video analysis application, I demonstrated that the 150 gendered pose prototypes derived from static images can serve as a robust analytical lens for classifying poses in video frames.

There are two key contributions of this work. First, the tool provides a ready-to-use, open-source platform that enables researchers to analyze gendered body posing in video without requiring GPU infrastructure or advanced technical expertise. Second, the custom mapping layer between MediaPipe's 33-landmark output and the original 17-keypoint OpenPose format demonstrates that cross-framework pose harmonization is feasible in practice.

From a sociological perspective, the temporal dimension introduced by video analysis offers a key advantage over static image analysis: the ability to trace gender performance as it unfolds over time. Preliminary testing of the tool indicates that individuals do not maintain a single ``gendered'' posture throughout a video; instead, they move between neutral and gender-coded clusters.




\section*{Statement on the Use of Artificial Intelligence Tools}

Artificial intelligence tools were used for code implementation, debugging, and for grammar and style suggestions during drafting this report. All substantive content, analysis, interpretation, and conclusions are the authors' own, and all sources were verified and cited by the authors. No AI tools were used to generate data or produce results. The authors reviewed and edited all AI-assisted text and code for accuracy and compliance with academic integrity policies.

% =============================================================================
% REFERENCES
% =============================================================================
\bibliographystyle{apalike}
\bibliography{references}


% =============================================================================
% APPENDIX
% =============================================================================
\newpage
\appendix
\section*{Appendix}
\addcontentsline{toc}{section}{Appendix}

\subsection*{A. System Pipeline}

% --- Uncomment once pipeline_diagram.png is available ---
% \begin{figure}[H]
%     \centering
%     \includegraphics[width=0.85\textwidth]{pipeline_diagram.png}
%     \caption{The computational pipeline: Video frames are processed via MediaPipe, normalized to the OpenPose format, and matched to the nearest of 150 prototypes.}
%     \label{fig:pipeline}
% \end{figure}

\noindent\textit{[Insert pipeline\_diagram.png here. The diagram should illustrate: Video $\rightarrow$ Frame extraction (0.5s intervals) $\rightarrow$ MediaPipe Pose (33 landmarks) $\rightarrow$ Mapping to 17-point format $\rightarrow$ Normalization (bounding box $\rightarrow$ origin $\rightarrow$ unit vector) $\rightarrow$ Euclidean distance to $>15,000$ dataset poses $\rightarrow$ Nearest pose's cluster assignment $\rightarrow$ Gender classification.]}

\subsection*{B. Keypoint Mapping}

% --- Uncomment once keypoint mapping figure is available ---
% \begin{figure}[H]
%     \centering
%     \includegraphics[width=0.7\textwidth]{keypoint_mapping.png}
%     \caption{Keypoint mapping strategy. The `Neck' and `MidHip' points (highlighted) are derived from MediaPipe landmarks to match the topology of the training dataset.}
%     \label{fig:keypoint_map}
% \end{figure}

\noindent\textit{[Insert keypoint\_mapping.png here. The diagram should show MediaPipe's 33-point skeleton alongside the target 17-point skeleton, with arrows showing direct mappings and the derived Neck/MidHip midpoints highlighted in a distinct color.]}

\subsection*{C. Cluster Prototype Examples}

Below are selected cluster prototypes generated by the independent K-Means clustering. Each stick figure represents the centroid (arithmetic mean) of all poses assigned to that cluster. The color bar beneath each prototype indicates the gender composition of the cluster: orange segments represent female-labeled poses, teal segments represent male-labeled poses, and blue segments represent non-binary-labeled poses.

Figure~\ref{fig:masc_prototypes} shows representative masculine-coded cluster prototypes, characterized by wide, symmetric stances and expansive postures.

% --- INSERT male_example_prototypes.png HERE ---
\begin{figure}[H]
    \centering
    \includegraphics[width=\textwidth]{data/output/male_example_prototypes.png}
    \caption{Selected masculine-coded cluster prototypes. Stick figures show the centroid pose of each cluster, with color bars beneath indicating gender composition (orange = women, teal = men, blue = non-binary). These clusters are characterized by wide stances, symmetric postures, and expansive body positioning.}
    \label{fig:masc_prototypes}
\end{figure}

% --- INSERT YOUR FEMININE PROTOTYPE VISUALIZATION HERE ---
% Replace the path below with the actual filename of your feminine prototypes figure.
\begin{figure}[H]
    \centering
    \includegraphics[width=\textwidth]{data/output/female_example_prototypes.png}
    \caption{Selected feminine-coded cluster prototypes. These clusters are characterized by crossed legs, asymmetric postures, canted body positions, and S-shaped body displays. Color bars show gender composition (orange = women, teal = men, blue = non-binary).}
    \label{fig:fem_prototypes}
\end{figure}

% --- INSERT YOUR NEUTRAL PROTOTYPE VISUALIZATION HERE ---
% Replace the path below with the actual filename of your neutral prototypes figure.
\begin{figure}[H]
    \centering
    \includegraphics[width=\textwidth]{data/output/neutral_example_prototypes.png}
    \caption{Selected neutral cluster prototypes showing no statistically significant gender imbalance. These clusters represent common poses adopted equally by all genders. Color bars show gender composition (orange = women, teal = men, blue = non-binary).}
    \label{fig:neutral_prototypes}
\end{figure}

\end{document}
